% !TeX root = booklet-02.tex
\ifCompleteNotes
\section*{راهنمای بخش دوم، اثرهای غیرایده‌آل}
\addcontentsline{toc}{section}{راهنمای بخش دوم، اثرهای غیرایده‌آل}
\else
\chapter*{راهنمای مطالعه}
\addcontentsline{toc}{chapter}{راهنمای مطالعه}
\fi
این نسخه از جزوه دوم، \lr{Mobility Degradation}، \lr{Velocity Saturation}،
\lr{Channel Length Modulation}، اثرهای ولتاژ آستانه، نشتی، دما و هندسه
را بر اساس بخش‌های \secnum{۲.۴.۱} تا \secnum{۲.۴.۷} کتاب آموزش می‌دهد.
مبنای تطبیق، صفحات چاپی ۷۴ تا ۸۷ کتاب، اسلایدهای ۳ تا ۳۷
\lr{Nonideal Transistor Theory} و تمرین‌های \secnum{۲.۲}، \secnum{۲.۳}، \secnum{۲.۶} و
\secnum{۲.۷} تا \secnum{۲.۱۰} است
\cite{weste,nonideal}. مطالعه قبلی \lr{Digital Electronics} فرض نشده است.
مرور پیش‌نیاز، معنای نمادها و محاسبات لازم در همین متن آمده‌اند.

ابتدا مرور کوتاه را بخوانید، سپس مسیر میدان، سرعت، بار و جریان را دنبال کنید.
پس از استخراج هر مدل، مثال عددی و کنترل واحدها ارائه می‌شود. روابط، شکل‌ها و
مثال‌های با عنوان «آموزشی» برای توضیح مدل تنظیم شده‌اند. نمودارهای جزوه از
روابط نوشته‌شده رسم شده‌اند و محاسبات آموزشی این مدل‌ها را نمایش می‌دهند.

تمرین \secnum{۲.۶} در متن استخراج مدل سرعت و تمرین‌های \secnum{۲.۳} و
\secnum{۲.۷} تا \secnum{۲.۹} در فصل
تمرین‌های اثر آستانه آمده‌اند. تمرین \secnum{۲.۱۰} پس از آموزش نشتی و پیش‌نیاز
نمایی و لگاریتم حل شده است. نتیجه تمرین \secnum{۲.۲} در مقدمه حل \secnum{۲.۳} استخراج
می‌شود. بحث دما، هندسه و جمع‌بندی، پیش از فصل‌های تمرین قرار دارد.
قراردادهای \lr{PVT} اسلایدهای ۳۱ تا ۳۴ نیز در یک مکمل توضیح داده شده‌اند.
مقادیر فناوری و آمار نشتی مرجع با زمینه زمانی نسخه سال ۲۰۱۱ بیان می‌شوند.

در تحلیل \lr{nMOS}، ولتاژ آستانه را مثبت می‌گیریم و قراردادهای زیر را
به کار می‌بریم.
\[
 V_{GT}=V_{GS}-V_t,\qquad V_t>0,\quad V_{DS}\ge0.
\]
نماد \(V_{GT}\) اضافه‌ولتاژ گیت نسبت به آستانه را نشان می‌دهد.

در تحلیل \lr{pMOS}، از ولتاژهای مثبت \(V_{SG}\) و \(V_{SD}\) و بزرگی
جریان استفاده می‌کنیم. بزرگی ولتاژ آستانه را به صورت زیر تعریف می‌کنیم.
\[
 V_{TP}=|V_{tp}|.
\]
در این تعریف، \(V_{tp}\) ولتاژ آستانه علامت‌دار \lr{pMOS} است.
نمادگذاری پایانه‌ها با کتاب متناظر است. برای نمونه، \(V_{GS}\) در جزوه
و \(V_{gs}\) در کتاب به یک کمیت اشاره دارند.

